\begin{claim}{A1.1}{Proved}{Local freeness and the determinant line}
Both $\mathcal A$ and $\mathcal E$ are locally free, of ranks three and
two.  The vertical vector $\xi=\partial/\partial\zeta_2$ is nowhere
zero as a section of $\mathcal E^\vee$.  There is an exact sequence
of holomorphic vector bundles
\begin{equation}\label{eq:A-E-line}
0\longrightarrow N\longrightarrow\mathcal E
 \xrightarrow{\langle\cdot,\xi\rangle}\OO_X\longrightarrow0,
\qquad N=\det\mathcal E=\omega_f(-2S_1-3S_2).
\end{equation}
\end{claim}
\begin{proof}
At a smooth point of a reduced cusp branch the assertion is the usual
coordinate calculation.  At a double point write $t_c=xy$ with third
coordinate $z$; then
$\mathcal A=R\,d\log x\oplus R\,d\log y\oplus R\,dz$ and its
base subline is $R(d\log x+d\log y)$.  At either triple point write
$t_c=xyz$; then $\mathcal A$ has the three logarithmic coordinate
basis vectors and its base subline is their sum.  In each case that
subline is a direct summand.  These calculations apply on every one
of the three double curves, including their endpoints; no simplicity
of the global divisor is assumed.

At $S_j$ we use the actual free cyclic quotient
$q:J'_j\to N_j$, which is finite and etale on total spaces, and
$t_j=s^{m_j}$ on the base.  On $J'_j$ the logarithmic conormal is
$m_j\,ds/s$, and its coefficient is a unit over $\C$.  It is again a
direct summand of the logarithmic cotangent bundle.  The comparison
$q^*\mathcal A=\Omega_{J'_j}^1(\log(s=0))$ identifies the quotient
bundles and their exact sequence; faithful flat descent proves the
claim downstairs.  This is not a replacement of $m_jS_j$ by $S_j$.

On the marked torus family $\xi$ is the constant second vertical
vector.  The linear monodromies fix it, and translations preserve it.
On a cusp toric chart it is a nonzero constant vector in the relative
logarithmic tangent lattice of rank two.  A toric coordinate change
$x_i\mapsto c_i(t_c)t_c^{b_i}x_i$ changes $d\log x_i$ by a multiple
of $d\log t_c$, so this constant relative vector descends and is
nowhere zero in that bundle.  It can vanish as an ordinary tangent
vector at a toric stratum; this does not imply vanishing in the
logarithmic bundle.  Hence evaluation is surjective and its kernel
is the determinant line.  Taking determinants of
$0\to f^*\mathcal L\to\mathcal A\to\mathcal E\to0$ gives
\[
\det\mathcal E=K_X(D)\otimes f^*\mathcal L^{-1}
=\omega_f\bigl(D-f^*\Sigma\bigr)
=\omega_f(-2S_1-3S_2),
\]
which proves \eqref{eq:A-E-line}.
\end{proof}
